\begin{table}[ht]
\centering
\caption{Contrastes planifiés sur les réponses perceptives « 2 flashs ».}
\label{tab:perceptual_response_contrasts_v2}
\scriptsize
\begin{tabular}{p{5.2cm}rrrrr}
\toprule
Contraste & Odds ratio & SE & $z$ & IC 95\% & $p$ \\
\midrule
A1V1+A2 standard vs A1V1 & 8.52 & 0.14 & 15.06 & [6.45, 11.26] & $< .001$ \\
A1V1+A2 court vs A1V1 & 15.17 & 0.14 & 19.33 & [11.51, 19.98] & $< .001$ \\
A1V1+V2 standard vs A1V1+A2V2 standard & 0.08 & 0.32 & -7.68 & [0.04, 0.16] & $< .001$ \\
A1V1+V2 court vs A1V1+A2V2 court & 0.17 & 0.14 & -13.26 & [0.13, 0.22] & $< .001$ \\
A1V1+A2 court vs A1V1+A2 standard & 1.78 & 0.12 & 4.94 & [1.42, 2.24] & $< .001$ \\
A1V1+V2 court vs A1V1+V2 standard & 0.17 & 0.14 & -13.26 & [0.13, 0.22] & $< .001$ \\
\bottomrule
\end{tabular}

\vspace{0.2cm}
\begin{minipage}{0.94\textwidth}
\footnotesize
\textit{Note.} Le modèle est un GLMM binomial prédisant la probabilité de répondre « 2 flashs » à partir de la condition expérimentale, avec un intercept aléatoire par participant : \texttt{response2 \textasciitilde{} condition\_cell + (1 | participant)}. Les contrastes sont exprimés en odds ratios. Une valeur supérieure à 1 indique une probabilité plus élevée de réponse « 2 flashs » dans la première condition du contraste.
\end{minipage}
\end{table}
